\title{Byte Latent Transformer: Patches Scale Better Than Tokens}

\begin{abstract}
We present the Byte Latent Transformer (\textsc{BLT}), a novel large language model architecture that operates directly at the byte level and, for the first time, achieves parity with traditional tokenization-based models at scale, while offering notable advances in inference speed and robustness. The \textsc{BLT} framework transforms input bytes into variable-length patches, utilizing these as fundamental computational entities. Patch boundaries are determined according to the entropy of the subsequent byte, enabling the allocation of computational and model resources to segments of higher data complexity as needed. We conduct the inaugural scaling study of byte-level architectures under controlled floating point operation constraints, scaling up to 8B model parameters and 4T bytes of training data. Our experiments affirm that it is possible to scale models trained exclusively on raw byte streams without relying on pre-defined vocabularies. Dynamic patch selection for predictable content not only enhances both training and inference throughput but also yields qualitative gains in reasoning and out-of-distribution generalization. Collectively, at equal inference resource budgets, \textsc{BLT} achieves superior scaling behavior compared to tokenization-based approaches, due to its simultaneous increase of both patch size and model complexity.
\end{abstract}

\section{Introduction}
\label{section:intro}

We present the Byte Latent Transformer~(\textbf{BLT}), a novel model architecture that eliminates the need for tokenization by directly processing sequences of raw bytes. BLT is the first approach of its kind to demonstrate competitive large-scale performance with traditional token-based methods, while also offering notable gains in both efficiency and resilience (\S\ref{sec:robustness}). In contrast, standard large language models (llms) are predominantly trained using end-to-end approaches, but still rely on a non-trainable tokenization step as a form of preliminary preprocessing. This step segments byte sequences into a fixed vocabulary of tokens, unintentionally encoding certain inductive biases in data compression. The reliance on tokenization introduces vulnerabilities, including heightened sensitivity to the specific domain or modality~\citep{dagangetting}, increased susceptibility to input perturbations~(\S\ref{sec:robustness}), deficiencies in representing orthographic distinctions~\citep{edman2024cute}, and unfair treatment of diverse languages~\citep{liang2023xlm,petrov2024language,limisiewicz2024myte}.

Tokenization has historically played a crucial role, as training large language models directly on bytes incurs substantial computational expense at scale due to extended sequence lengths~\citep{xue2022byt5}. Prior research has sought to address this by leveraging more computationally efficient self-attention mechanisms~\citep{el2020characterbert, clark2022canine} or utilizing architectures that eliminate attention entirely~\citep{wang2024mambabyte}~(\S\ref{section:related_work}). Nevertheless, such strategies mainly benefit the training of \textit{small models}. When models are scaled up, the primary computational bottleneck becomes the extensive feed-forward network layers—applied to each byte—rather than the cost associated with the attention mechanism itself.

To optimize compute allocation, we introduce an adaptive and learnable approach for aggregating bytes into \textit{patches}~(\S\ref{section:patching}), alongside a novel architecture that integrates both byte-level and patch-level signals. In contrast to tokenization, {\textsc BLT} does not rely on a predetermined patch vocabulary. Instead, flexible combinations of bytes are projected into latent patch embeddings through compact, trainable encoder and decoder components. Our findings indicate that this strategy allows for \textit{greater} efficiency in compute distribution relative to models built around tokenization.

Tokenization-based LLMs assign an equal amount of computational effort to each token. This approach prioritizes uniformity over efficiency, as the heuristics used for compression in token induction do not consistently align with the inherent difficulty of making accurate predictions. Our approach centers around the principle that computational resources should be flexibly distributed according to task demands. For instance, predicting the end of most words is a relatively straightforward, low-entropy task that does not require the capacity of a large transformer, in contrast to selecting the initial token in a sentence, which involves greater complexity. This notion informs the design of {\textsc BLT}'s architecture~(\S\ref{section:architecture}), which utilizes three transformer components: a pair of compact byte-level \textit{local models} complemented by a large, global \textit{latent transformer}~(\autoref{fig:arch_high_level}). 

To organize bytes into appropriate patches and enable adaptive computation, {\textsc BLT} employs segmentation of input based on the entropy observed in predicting subsequent bytes. This results in context-aware groupings where information density is relatively homogeneous.

We introduce the first comprehensive study on flop-controlled scaling of byte-level models, extending up to 8 billion parameters and 4 trillion bytes of training data. Our findings demonstrate that it is feasible to train large-scale, end-to-end models directly from raw bytes, thereby obviating the need for fixed-vocabulary tokenization. Across the board, {\textsc BLT} achieves training performance, as measured by computational cost in flops,\footnote{We define a model's computational requirement by enumerating the total Floating Point OPerations (flops) performed.} comparable to that of Llama 3, yet with as much as 50\% reduction in inference flops~(\S\ref{section:scaling}).

Additionally, our results reveal that handling raw byte inputs significantly enhances the modeling of infrequent or atypical data. {\textsc BLT} models not only show greater resilience to input noise than tokenizer-dependent counterparts but also exhibit improved performance at the character level, excelling in tasks such as orthography, phonological processing, and low-resource machine translation~(\S\ref{sec:robustness}).

Moreover, the {\textsc BLT} architecture allows for simultaneous increases in both model capacity and patch size while adhering to a fixed inference flop constraint. Using longer patch sizes can yield computational savings during inference, which in turn can be allocated to scale the global latent transformer, since it is invoked less frequently. Through scaling experiments controlled for inference flops~(\autoref{fig:fixed_inference_scaling}), we observe more favorable scaling patterns compared to models relying on token-based preprocessing.

To conclude, this work presents several key contributions: 1) We propose {\textsc BLT}, a byte latent LLM framework that adaptively distributes computational resources to enhance flop efficiency; 2) We demonstrate that our approach matches the training flop efficiency of Llama 3 at scales up to 8B parameters, and allows a tradeoff between modest decreases in evaluation scores and substantial improvements in flop efficiency, reaching gains as high as 50\%; 3) {\textsc BLT} introduces a novel pathway for scaling LLMs, permitting increases in model capacity without exceeding a predetermined inference budget; 4) Our experiments reveal that {\textsc BLT} exhibits enhanced robustness to input perturbations and possesses greater sensitivity to sub-word elements of input sequences, outperforming traditional token-based LLMs in this regard. The corresponding code for training and inference is made available at~\url{https://github.com/facebookresearch/blt}.

\section{Patching: From Individual Bytes to Groups of Bytes}
\label{section:patching}

Dividing byte sequences into \textit{patches} enables {\textsc BLT} to flexibly distribute computational resources depending on contextual requirements. Several approaches to segmenting bytes into patches are depicted in Figure~\ref{fig:patching_types}. More formally, a patching function $f_p$ partitions a byte sequence $\pmb{x}=\{x_i,|i=1,\ldots n\}$ of length $n$ into a series of $m < n$ patches $\pmb{p}=\{p_j|j=1,\ldots,m\}$, by associating each $x_i$ with \{0,1\}, where 1 marks the beginning of a new patch. For both token-centric and patch-based architectures, the major determinant of computational expenditure arises from the total number of operations performed by the main Transformer, which in {\textsc BLT} corresponds to the patch count required to represent the data as per the chosen patching function. Therefore, the mean patch length, or simply \textit{patch size}, serves as the principal determinant of computational costs encountered during both training and inference for a particular patching method~(\S\ref{section:flops}).

We proceed by introducing three strategies for constructing patches: using a fixed number of bytes per patch~(\S\ref{section:static-patch}), leveraging whitespace boundaries~(\S\ref{section:white-patch}), and dynamically segmenting based on entropy values from a compact byte language model~(\S\ref{section:dyn-patch}). To conclude, we consider incremental patching strategies as well as distinguish between patching and standard tokenization~(\S\ref{section:bpe-patch}).

\subsection{Strided Patching Every K Bytes}
\label{section:static-patch}

A common approach for grouping bytes involves forming patches of a fixed length $k$, as employed in MegaByte~\citep{yu2023megabyte}. This method utilizes a constant stride, which simplifies both the training and inference processes, offers a convenient means of adjusting the average patch size, and allows for straightforward manipulation of computational cost. Despite these advantages, the technique has notable drawbacks. Primarily, computation is assigned uniformly, disregarding areas where more or less processing might be required; for example, transformer step $j$ may be squandered on predicting whitespace in code, or fail to provide adequate computation when encountering information-rich regions such as mathematical text. Additionally, this static grouping can cause inconsistent and context-insensitive segmentations of equivalent byte sequences, resulting in, for instance, the arbitrary splitting of identical words.

\subsection{Space Patching}
\label{section:white-patch}

\citet{slagle2024spacebyte} introduces a straightforward yet powerful modification to strided patching that generates new patches following any space-like bytes\footnote{
    Space-like bytes refer to any byte except for Latin characters, digits, or UTF-8 continuation bytes. Furthermore, each patch must include at least one non space-like byte.}, which often correspond to natural linguistic unit boundaries across languages. In Space patching, an additional latent transformer operation (incurring greater computational cost) is dedicated to the modeling of each word. This method ensures that individual words are consistently patched within sequences, directing computational resources (flops) to the positions with high uncertainty, such as just after spaces. As an illustration, generating the initial byte for the answer to “Who composed the Magic Flute?\uline{\hspace{.6em}}” presents a greater challenge than predicting subsequent bytes after ``M," because the leading letter dramatically narrows the prediction space, thereby making the remainder of the completion (``ozart'') much more straightforward. Despite these benefits, space patching does not effectively support all domains and languages, and more critically, is incapable of adjusting patch sizes. In the following section, we propose an alternative patching technique, leveraging the observation that the initial bytes of words usually present the greatest prediction difficulty, while also enabling flexible patch sizing.

\subsection{Entropy Patching: Using Next-Byte Entropies from a Small Byte LM}
\label{section:dyn-patch}

Instead of employing heuristic rules like splitting on whitespace, we adopt a data-centric strategy to detect locations where next-byte predictions exhibit high uncertainty. Our proposed method, \textit{entropy patching}, leverages entropy estimates to determine the locations of patch boundaries.

We train a small byte-level auto-regressive language model on the training data for {\textsc BLT} and compute next byte entropies under the LM distribution $p_{e}$ over the byte vocabulary $\mathcal{V}$:

\begin{equation}
H(x_i) = \sum_{v \in \mathcal{V}} p_{e}(x_i=v|\pmb{x}_{<i}) \log p_{e}(x_i=v|\pmb{x}_{<i})
\end{equation}

We investigate two approaches for detecting patch boundaries using entropies $H(x_i)$. The first method selects locations where the entropy surpasses a universal threshold, as depicted in~\autoref{fig:patching}. The second method highlights points whose entropy values are significantly greater than the immediately preceding entropy. This latter technique may also be seen as detecting points where the approximate monotonic decline in entropy within a patch is interrupted.

\begin{align*}
    \text{Global Constraint}\hspace{1cm}H(x_t)& > \theta_g\\
    \text{Approx. Monotonic Constraint}\hspace{1cm}H(x_t)& - H(x_{t-1}) > \theta_r
\end{align*}

During the data loading phase, patch boundaries are detected through an efficient preprocessing procedure. This contrasts with the approach of~\citet{nawrot-etal-2023-efficient}, who employ a classifier trained to infer patch boundaries based on entropy measures. In our empirical evaluation~(\S\ref{section:experiments}), we examine and contrast these two strategies for identifying bytes of low versus high entropy.

\subsection{The Byte-Pair Encoding (BPE) Tokenizer and Incremental Patching}
\label{section:incremental-patching}
\label{section:bpe-patch}

Numerous contemporary llms, such as the baseline Llama 3, employ subword tokenization methods like bpe~\citep{gage1994new, sennrich-etal-2016-neural}. Throughout this work, we use the term ``tokens'' to denote groups of bytes selected from a \textit{finite} vocabulary established before model training, contrasting them with ``patches,'' which represent adaptively formed sequences that are not restricted by a predetermined vocabulary. Notably, tokens differ fundamentally from patches in that they prevent the model from directly leveraging raw byte-level information.

An important advancement of {\textsc BLT} relative to models reliant on tokenization is its novel approach to balancing vocabulary size against computational requirements. Traditional llms face a compromise: expanding the vocabulary allows for larger average tokens, which reduces the number of processing steps required per sequence, but also necessitates a higher-dimensional output in the model's final projection layer. This compromise significantly restricts the flexibility of tokenization-based strategies in altering token length or inference overhead. As an illustration, Llama 3 achieves an increase in average token size from 3.7 to 4.4 bytes, but this comes at the expense of quadrupling its embedding table’s capacity compared to Llama 2.

During generation, {\textsc BLT} must determine at each step whether the current position in the byte sequence marks the beginning of a new patch, since this choice dictates whether to activate additional computation through the Latent Transformer. Importantly, this decision must be made without knowledge of upcoming bytes, as the rest of the sequence has not yet been produced. Therefore, the patching process must segment the byte sequence based solely on past and present information, not on the unseen future. More formally, a patching scheme $f_p$ exhibits the property of incremental patching if it holds that:
$$f_p(\pmb{x}_{<i}) = f_p(\pmb{x})_{<i}$$
bpe fails to qualify as an incremental patching scheme, because the same initial segment may yield different tokenizations contingent upon the subsequent content of the sequence, violating the criterion above\footnote{Employing a dedicated delimiter token to signal patch boundaries would enable bpe to conform to an incremental patching approach; however, this increases the total byte-sequence length.}.

\section{{\textsc BLT} Architecture}
\label{section:architecture}
The {\textsc BLT} framework integrates a large global autoregressive language model designed to process patch representations, paired with two compact local modules. These local components are responsible for transforming byte sequences into patches and, conversely, reconstructing byte sequences from patch representations~(\autoref{fig:arch_high_level}).

\subsection{Latent Global Transformer Model}

\textit{The Latent Global Transformer} refers to an autoregressive transformer architecture $\mathcal{G}$ with $l_{\mathcal{G}}$ layers. This model transforms a series of latent input patch representations, $p_j$, into a corresponding sequence of output patch representations, $o_j$. In this work, the index $j$ designates different patches, while $i$ denotes individual bytes. Employing a block-causal attention mask~\citep{dubey2024llama}, the global model ensures that attention is permitted only up to and including the current patch within each document. As most of the computational cost in both pre-training and inference is attributed to this model, strategically deciding when to execute $\mathcal{G}$ enables selective allocation of computational resources across input and output segments, adapting to their specific complexities.

\subsection{Local Encoder}
\label{section:local}

The \textit{Local Encoder Model}, represented as $\mathcal{E}$, utilizes a compact transformer architecture comprising $l_{\mathcal{E}} << l_{\mathcal{G}}$ layers. Its principal objective is the rapid transformation of an input byte sequence $b_i$ into rich patch representations $p_j$. Distinct from conventional transformer setups, each transformer layer in $\mathcal{E}$ is succeeded by a cross-attention mechanism, which consolidates byte-level features into patch-level representations, as illustrated in~(\autoref{fig:crossattn}).

The bytes $b_i$ are first projected into continuous space via embeddings using a matrix of shape $\mathbb{R}^{256 \times h_{\mathcal{E}}}$, resulting in vectors $x_i$. These initial embeddings may be further enhanced by hash-embeddings as outlined in~(\S\ref{sec:hash-emb}).

Subsequent layers alternate between transformer and cross-attention modules~(\S\ref{sec:enc-cross-attn}), iteratively refining the encoded representations into patch vectors $p_i$, which are ultimately input to the global transformer model $\mathcal{G}$. The attention mask employed within the transformer layers is of the \textit{local block causal} variant: each position attends to a set window of $w_{\mathcal{E}}$ previous bytes. While these windows can span across dynamic patch boundaries, they strictly do not breach document boundaries. The next subsections elaborate on specifics of the embedding mechanisms and the structure of the cross-attention layer.

\subsubsection{Encoder Hash n-gram Embeddings}
\label{sec:ngram-emb}
\label{sec:hash-emb}
An essential aspect of generating strong and informative representations at any given position $i$ lies in leveraging knowledge of prior bytes. In {\textsc BLT}, this is addressed by encoding not only the isolated byte $b_i$ but also embedding it within a sequence of bytes, i.e., a byte n-gram. Specifically, for each position $i$, we initially generate byte-grams

\begin{align}
    g_{i,n}=\{b_{i-n + 1},\ldots, b_i\}
\end{align}

for each byte position $i$ and for $n$ ranging from three to eight.\footnote{
Byte-grams of length $n$ or greater are disregarded in cases where $i<n$.}

Subsequently, we propose hash $n$-gram embeddings, whereby every byte $n$-gram is projected through a hash function onto an index within a designated embedding matrix $E_{n}^{hash}$ of fixed dimension, for every $n$ belonging to $\{3,4,5,6,7,8\}$~\citep{bai2010learning}. 
This generated embedding is then combined with the original byte embedding, and the sum is normalized prior to serving as the input to the local encoder module. The formulation for the enhanced embedding is as follows:

\begin{align}
e_i &= x_i + \sum_{n = 3,...,8}E_{n}^{hash}(\text{Hash}(g_{i,n})) \\
\text{where, Hash}(g_{i,n}) &= \text{RollPolyHash}(g_{i,n}) \% |E_{n}^{hash}|
\end{align}

We scale $e_i$ by dividing it by the number of $n$-gram sizes incremented by one, employing the RollPolyHash method as specified in Appendix~\ref{appendix:rollpolyhash}. Section~\ref{section:ablations} explores ablation studies on the impact of $n$-gram hash embeddings across varying $n$ values and embedding table sizes, examining their influence on flop-controlled scaling law behaviors. Beyond hash-based $n$-gram embeddings, we also conducted experiments using frequency-based $n$-gram embeddings; further information on this investigation is provided in Appendix~\ref{appendix:ngrams}.

\subsubsection{Encoder Multi-Headed Cross-Attention}
\label{sec:enc-cross-attn}
Our approach builds directly upon the cross-attention mechanism employed in the Perceiver architecture~\citep{jaegle2021perceiver}, with the primary modification being that the latent variables now represent patch-specific encodings, rather than being drawn from a static pool of latent representations~(\autoref{fig:crossattn}); additionally, each patch is designed to attend exclusively to its associated byte sequence. For each patch $p_j$, the attention module introduces a query vector, which is constructed by first aggregating the corresponding byte-level embeddings for patch $p_j$ and then subjecting this pooled representation to a linear transformation, denoted $\mathcal{E}_{C} \in \mathbb{R}^{h_{\mathcal{E}} \times (h_{\mathcal{E}} \times U_{\mathcal{E}})}$, with $U_{\mathcal{E}}$ representing the encoder’s cross-attention head count. Concretely, letting $f_{\text{bytes}}(p_j)$ indicate the byte sequence assigned to patch $p_j$, the computation proceeds as follows:

\begin{align}
P_{0,j} &= \mathcal{E}_{C}(f_\text{bytes}((p_j)), f~ \text{is a pooling function} \\
P_l &= P_{l-1} + W_o\left(\text{softmax}\left(\frac{QK^T}{\sqrt{d_k}}\right)V\right) \\
\text{where } Q_j &= W_q(P_{l-1,j}), K_i = W_k(h_{l-1,i}), V_i = W_v(h_{l-1,i}) \\
h_l &= \text{Encoder-Transformer-Layer}_l(h_{l-1})
\end{align}

Here, $P \in \mathbb{R}^{n_p \times h_{\mathcal{G}}}$ denotes $n_p$ patch-level representations, which are intended for input into the global model. These patch representations are initially constructed by aggregating the byte embeddings $e_i$ associated with each patch $p_j$. The matrices $W_q$, $W_k$, $W_v$, and $W_o$ act as the linear transformations for queries, keys, values, and outputs, respectively. Specifically, the keys and values are obtained by applying these projections to the byte representations $h_i$ from the previous layer (using $e_i$ for the initial layer). To ensure attention operates within patch boundaries, we implement a targeted masking scheme such that every query $Q_j$ only interacts with keys and values tied to the bytes within its own patch $j$. Employing multi-head attention across $Q, K$, and $V$, and because patch vectors generally have a greater dimensionality ($h_{\mathcal{G}}$) compared to $h_{\mathcal{E}}$, we organize $P_l$ as a set of multi-head vectors with each head sized $h_{\mathcal{E}}$ during the cross-attention step, subsequently concatenating these to achieve the full $h_{\mathcal{G}}$ dimensional representation. Furthermore, we apply pre-LayerNorm normalization to the queries, keys, and values prior to attention computation, and do not incorporate positional embeddings in this cross-attention mechanism. The cross-attention block itself is surrounded by a residual connection.

\subsection{Local Decoder} 

The local decoder $\mathcal{D}$, paralleling the structure of the local encoder, is implemented as a compact transformer architecture composed of $l_{\mathcal{D}}$ layers, with $l_{\mathcal{D}} << l_{\mathcal{G}}$. Its role is to transform the sequence of global patch representations $o_j$ back into the sequence of raw bytes $y_i$. In generating raw byte predictions, the decoder conditions on previously reconstructed bytes, utilizing the latent features obtained from the local encoder for the corresponding byte sequence. The decoding process involves $l_{\mathcal{D}}$ successive layers, alternating between cross-attention and standard transformer operations. Each cross-attention mechanism is performed prior to the transformer layer in order to synthesize byte-level representations from the patch-level inputs, after which the local decoder’s transformer layer further processes these byte representations.

\subsubsection{Decoder Multi-headed Cross-Attention} 

In the decoder cross-attention, the roles of the queries and key/values are interchanged i.e. the byte-representations are now the queries, and the patch representations are now the key/values. The initial byte-representations for the cross-attention are initialized as the byte embeddings from the last encoder layer i.e. $h_{l_{\mathcal{E}}}$. The subsequent byte-representations for layer $l$, $d_{l,i}$ are computed as:

\begin{align}
D_0 &= h_{l_{\mathcal{E}}} \\
B_l &= D_{l-1} + W_o\left(\text{softmax}\left(\frac{QK^T}{\sqrt{d_k}}\right)V\right), \\
  &\text{where } Q_i = W_q(d_{l-1,i}), K_i = W_k(\mathcal{D}_C(o_j)), V_i = W_v(\mathcal{D}_C(o_j)) \\
  D_l &= \text{Decoder-Transformer-layer}_l(B_l)
\end{align}

Here, $W_k$ and $W_v$ denote the key and value projection matrices, respectively, both of which carry out linear transformations followed by the split operation $\mathcal{D}_C$ on the final patch representations $o_j$ produced by the global model. The query projection matrix $W_q$ processes byte-level representations $d_{l-1}$ from the preceding decoder transformer layer (or $h_{l_{\mathcal{E}}}$ for the initial layer). The output is subsequently mapped through the output projection matrix $W_o$, producing $B \in \mathbb{R}^{h_{\mathcal{D}} \times n_b}$, where $n_b$ designates the number of output bytes. The resulting decoder representations, $D_l$, are obtained by feeding the output $B$ of the cross-attention module into a decoder transformer layer. Similar to the configuration in the local encoder’s cross-attention, our architecture employs multi-head attention, incorporates pre-layer normalization, omits positional embeddings, and applies a residual connection around the cross-attention block.

\section{Experimental Setup}
\label{section:experiments}
We implemented rigorously controlled experimental conditions to fairly evaluate {\textsc BLT} against models relying on traditional tokenization, ensuring that {\textsc BLT} did not benefit from increased context window lengths relative to the baselines.

\subsection{Pre-training Datasets} In this work, all model sizes are initially pre-trained on two distinct datasets: 1) The Llama 2 dataset~\citep{touvron2023llama}, consisting of 2 trillion tokens sourced from a wide range of publicly accessible data, followed by extensive cleaning and filtering processes to enhance overall quality; and 2) {\textsc BLT}-1T: a newly constructed dataset containing 1 trillion tokens from diverse public datasets, supplemented by a subset of the pre-training corpus released as part of Datacomp-LM~\citep{li2024datacomp}. The Llama 2 dataset is specifically leveraged to conduct scaling law analyses, employing the token counts found optimal by~\cite{dubey2024llama} in order to inform architectural decisions for {\textsc BLT}. In contrast, the {\textsc BLT}-1T dataset is utilized for comprehensive pre-training, facilitating a direct comparison against Llama 3 in downstream evaluations. Notably, both datasets exclude any data originating from Meta-affiliated platforms or services. For all baseline comparisons involving models based on tokenization, we adopt the Llama 3 tokenizer, which features a vocabulary size of 128K tokens; this choice consistently yielded superior baseline results relative to the Llama 2 tokenizer within our experimental setup.

\subsection{Entropy Model}  
For our experiments, the entropy model consists of a byte-level language model, which is trained on the same data distribution as the complete {\textsc BLT} model. Except where specified, we employ a transformer architecture comprising 100 million parameters, organized into 14 layers with a hidden size of 512, and utilizing sliding window attention spanning 512 bytes. All other hyperparameter settings mirror those used in both our local and global transformers. We have explored variations in model scale, receptive field length, and architectural choices, as outlined in \autoref{section:ablations}. Notably, if the model's receptive field is sufficiently constrained, it is possible to represent the learned entropy model compactly as a lookup table.

\subsection{Entropy Threshold and Equalizing Context Length} For entropy-based patching approaches, we determine a threshold that yields a target mean \textit{patch size} across the pretraining dataset. In {\textsc BLT}, in contrast to standard tokenization, the selection of \textit{patch size} is flexible, substantially influencing the effective context window available to the model. To avoid conferring an unfair benefit to configurations with larger patch sizes, we fix the expected number of bytes per batch, thereby standardizing the average context length. Consequently, models employing larger patch sizes are trained with correspondingly shorter sequence lengths. Specifically, we adopt an 8k-byte context for Llama 2 data, whereas for the {\textsc BLT}-1T corpus, the average context is expanded to 16k bytes, with the batch size consistently set to 16M bytes.

Although the average batch size remains unchanged, dynamic patching approaches produce varying proportions of bytes per patch during data batching. To maximize efficiency, our implementation of {\textsc BLT} training aggregates patches in each batch such that padding within the resource-intensive latent transformer is unnecessary; as a result, each batch maintains a consistent patch count. For the purpose of training, we either pad or, if necessary, truncate input byte sequences to lengths of 12k bytes for the Llama 2 dataset and 24k bytes for the {\textsc BLT}-1T dataset. This procedure is designed to prevent memory usage surges that could arise from exceptionally large patches within some sequences.

\subsection{Entropy Model Context}
\label{section:entropy-context}
Our empirical results indicate that entropy patching tends to produce increasingly large patches in contexts featuring structured content, such as multiple choice tasks (refer to the patching behavior illustrated with an MMLU example in \autoref{fig:mmlu_patching}), where content repetitiveness is common. The observed variations stem from reduced entropy within regions of recurrent content in the entropy model context. Consequently, in our large-scale application of {\textsc BLT}-Entropy with a patch size of 4.5, we refresh the entropy context by introducing new lines and implement an approximate monotonicity constraint, since this approach is less vulnerable to "entropy drift" resulting from alterations in context length. Notably, this modification solely influences the way we compute entropies; the process used to determine the entropy threshold remains unchanged.

\subsection{FLOPs Estimation} 
\label{section:flops}

Our approach is primarily based on the methodology outlined in Chinchilla~\citep{hoffmann2022training} for calculating transformer flops, which accounts for operations in the feed-forward networks, the qkvo projections within the self-attention mechanism, as well as the computation required for attention and the output projection. However, a key distinction is that we treat the input embedding layer as an optimized lookup table rather than performing dense matrix multiplication, resulting in this layer contributing zero flops. Consistent with established literature, we assume that the backward pass incurs twice as many flops as the forward pass.

To compute flops \textit{per byte} for {\textsc BLT} models, we add up the flops for the local encoder transformer, the global latent transformer, and the local decoder transformer, together with the cross attention blocks in the encoder and the decoder:

\begin{align}
    \text{FL}_{\text{{\textsc BLT}}} &= \text{Transf. FL}(h_{\mathcal{G}}, l_{\mathcal{G}}, m=n_{ctx}/n_p,V=0) / n_p\\
   &+\text{Transf. FL}(h_{\mathcal{E}}, l_{\mathcal{E}}, m=w_{\mathcal{E}}, V=0) \\
   &+ \text{Transf. FL}(h_{\mathcal{D}}, l_{\mathcal{D}}, m=w_{\mathcal{D}}, V=256) \\ 
    &+ \text{Cross Attn. FL}(h_{\mathcal{E}}, l_{\mathcal{E}}, m=n_p, r=n_p/k) \cdot \frac{k}{n_p} \\ 
   &+ \text{Cross Attn. FL}(h_{\mathcal{D}}, l_{\mathcal{D}}, m=k, r=k/n_p) 
\end{align}

Here, $n_{ctx}$ denotes the sequence length measured in bytes, $n_p$ refers to the patch size, $r$ represents the ratio between query vectors and key/value vectors, and $k$ defines the relationship between the patch dimension and the byte dimension; specifically, it is the count of local model partitions merged to generate the overall model representation (for example, $k=2$ in \autoref{fig:crossattn}). The symbol $V$ indicates the vocabulary size utilized in the output projection, which is present only within the local decoder. The context length $m$ employed by the attention mechanism varies based on whether the computation is performed over a byte sequence or a patch sequence. The attention-related floating point operations (flops) are adjusted appropriately for each structural module. Precise formulations for calculating the flops in both the Transformer and Cross-Attention modules can be found in Appendix~\ref{section:flops-appendix}.

\subsection{Bits-Per-Byte Estimation} Perplexity only makes sense in the context of a fixed tokenizer as it is a measure of the uncertainty for each token. When comparing byte and token-level models, following previous work~\citep{xue2022byt5, yu2023megabyte, wang2024mambabyte}, we instead report Bits-Per-Byte (BPB), a tokenizer independent version of perplexity. Specifically:

\begin{align}
    \text{BPB}(x)&= \frac{\mathcal{L}_{CE}(\pmb{x})}{\ln(2)\cdot n_{\text{bytes}}}
\end{align}

Here, the aggregate cross-entropy loss, which quantifies the uncertainty associated with the data $\pmb{x}$, is divided by both the number of bytes present in $\pmb{x}$ and a fixed scaling factor to yield a normalized measure.

\subsection{Transformer Architecture Hyperparameters} In designing all transformer components within {\textsc BLT}, encompassing both local and global submodels, our configuration is principally based on the design choices established in Llama~3~\citep{dubey2024llama}. Specifically, we incorporate the SwiGLU activation function~\citep{swiglu} within the feed-forward modules, and employ rotary positional embeddings (RoPE)~\citep{RoPe} in all self-attention layers, parameterized by $\theta=500000$ as suggested by~\citep{xiong2024effective}. For normalization, RMSNorm~\citep{RMSNorm} is utilized across all layers. To optimize the computation of self-attention in scenarios utilizing deterministic attention masks—such as \textit{block causal} or \textit{fixed-window block causal}—we adopt Flash attention~\citep{dao2022flashattention}, and maintain a window span of 512 tokens for fixed-width masks. In contrast, cross-attention layers, which require adaptive, patch-specific masking, are accelerated via Flex Attention\footnote{\url{https://pytorch.org/blog/flexattention}}, enabling the use of highly optimized fused implementations that markedly enhance training throughput.

\subsection{{\textsc BLT}-Specific Hyperparameters} 

To evaluate the performance of the {\textsc BLT} models, we carry out analyses focusing on two main aspects: the models’ scaling properties and their effectiveness on downstream tasks. Our investigations encompass a range of model sizes—specifically 400M, 1B, 2B, 4B, and 8B parameters. Further architectural details for these models are provided in Appendix~Table~\ref{tab:arch}.

In the local encoder, query vectors for the initial cross-attention layer are generated using max-pooling. For all variants of the {\textsc BLT} model, we apply $500,000$ hashes using a single hash function and utilize n-gram lengths between 3 and 8. A uniform learning rate of $4e-4$ is adopted across all model sizes. This rate is selected because, in a hyperparameter search at the 400M and 1B model levels exploring the interval from $1e-3$ to $1e-4$, a consistent optimal rate was observed for both the token and {\textsc BLT} models.

In our scaling law experiments with Llama-2 datasets, the batch sizes correspond to those suggested in~\cite{dubey2024llama} or their byte-equivalent. Model training employs the AdamW optimizer~\citep{adamw}, using $\beta_1=0.9$, $\beta_2=0.95$, and $\epsilon=10^{-8}$. Learning rate scheduling involves 2000 steps of linear warm-up, followed by cosine decay down to zero. Throughout, a weight decay coefficient of 0.1 is maintained, and gradients are globally clipped at a maximum norm of 1.0.

\section{Scaling Trends}
\label{section:scaling}

We offer a comprehensive overview of the scaling dynamics observed in byte-level models, providing insights that may guide the future scaling of {\textsc BLT} models. Our investigation seeks to overcome the shortcomings of earlier studies on byte-level approaches through the following methods: (a) by analyzing scaling patterns within a compute-optimal training context, (b) by constructing equivalent 8B models with substantial training data exposure (reaching up to 1T tokens or 4T bytes) and systematically evaluating their performance on various downstream applications, and (c) by assessing scaling properties under fixed inference-cost conditions. Subsequent sections will explore in detail the particular benefits of utilizing byte-sequence modeling.

\subsection{Parameter Matched Compute Optimal Scaling Trends}
\label{section:parameter-matched-scaling}
Utilizing the Llama 2 dataset, we conduct training of several \textit{compute-optimal} bpe and {\textsc BLT} models, spanning four parameter scales between 1B and 8B. We subsequently chart the relationship between training flops and language modeling accuracy on a representative fraction of the combined training data. For the bpe models, the chosen model size relative to data quantity aligns with the compute-optimal proportion established in Llama~3~\citep{dubey2024llama}. This approach aims to optimize performance on the training set under a fixed computational budget, consistent with theoretical best practices outlined in~\citep{hoffmann2022training}, and serves as a strong reference point for our evaluations. Alongside each bpe model, a {\textsc BLT} counterpart is trained using the same input, wherein the Latent Transformer is configured to match both the architectural design and the scale of the paired bpe Transformer.

As shown in~\autoref{fig:scaling-compute-opt} (right), {\textsc BLT} models consistently achieve performance that is on par with or exceeds that of their bpe equivalents, and this pattern persists across increases in both model size and floating-point operations. To our knowledge, {\textsc BLT} represents the first Transformer architecture operating at the byte level to exhibit scaling behavior comparable to that of BPE-based models in compute-efficient settings. These findings thus support our hypothesis that the optimal parameter-to-compute ratio established for bpe holds for {\textsc BLT} as well, or, at the very least, deviates only minimally.

To align with bpe scaling trends, both advances in model architecture and the use of dynamic patching are essential. As depicted in ~\autoref{fig:scaling-compute-opt} (left), we assess models utilizing space-patching in relation to Llama 3. SpaceByte \citep{slagle2024spacebyte} is approximated here by employing {\textsc BLT} space-patching while omitting n-gram embeddings and cross-attention mechanisms. While SpaceByte offers improvements relative to Megabyte, its performance still lags behind that of Llama 3. \autoref{fig:scaling-compute-opt} (right) displays the performance gains contributed by both enhancements in architecture and the integration of dynamic patching. At scale, {\textsc BLT} models reach performance levels comparable to the most advanced tokenizer-based approaches like Llama 3.

Additionally, we note that the selection of tokenizer significantly influences model performance for tokenizer-based architectures; specifically, when trained on identical datasets, models utilizing the Llama-3 tokenizer exhibit superior results compared to those employing the Llama-2 tokenizer.

In summary, the {\textsc BLT} architecture demonstrates interpolation between the performance profiles of Llama 2 and Llama 3 when larger patch sizes are employed. While the bpe tokenizers used in Llama 2 and 3 exhibit average token sizes of 3.7 and 4.4 bytes, respectively, {\textsc BLT} achieves comparable scaling behaviors with average patch sizes of 6 and even up to 8 bytes. Because inference flops scale inversely with average patch size, selecting an 8-byte patch size can result in nearly a 50\% reduction in inference flops. Moreover, as both model and dataset sizes increase, models utilizing greater patch sizes appear to yield enhanced performance. Although {\textsc BLT} configured with an 8-byte patch size initially underperforms relative to bpe Llama 2 at the 1B parameter scale, it surpasses bpe at 7B. This trend implies that employing such large patch sizes may be even more advantageous at greater model and data scales, and that leveraging even larger patch sizes could become viable as models and training budgets continue to grow.

\subsection{Beyond Compute Optimal Task Evaluations}
\label{section:evals}
To further explore scaling characteristics, we train an 8B {\textsc BLT} model surpassing the compute-optimal threshold using the {\textsc BLT}-1T dataset—a larger, higher-quality corpus—and evaluate its effectiveness on a broad array of established classification and generation tasks. For the assessment, we focus on a range of benchmarks spanning common sense reasoning, factual knowledge, and code generation:
\paragraph{Classification tasks} consist of ARC-Easy (0-shot)~\citep{Eval_ARC}, ARC-Challenge (0-shot)~\citep{Eval_ARC}, HellaSwag (0-shot)~\citep{Eval_hellaswag}, PIQA (0-shot)~\citep{Eval_piqa}, and MMLU (5-shot)~\citep{hendrycks2020measuring}. For these, we utilize a prompt-scoring strategy based on the computed likelihood of candidate options, and present mean accuracy across datasets.
\paragraph{Coding related generation tasks:} Performance on MBPP (3-shot)~\citep{Eval_MBPP} and HumanEval (0-shot)~\citep{Eval_HumanEval} is measured using pass@1 scores, allowing for evaluation of large language model proficiency in producing correct Python code solutions.

In \autoref{tab:evals}, the performance of three distinct models trained on the {\textsc BLT}-1T dataset is presented: a Llama 3 tokenizer-based model that relies on bpe,\footnote{We opt for the Llama 3 tokenizer, which uses a 128k vocabulary, as it outperforms Llama 2’s 32k vocabulary.} alongside two configurations of the {\textsc BLT} model. The first variant incorporates a space-patching mechanism ({\textsc BLT}-Space), whereas the second adopts an entropy-based patching approach ({\textsc BLT}-Entropy). The latter applies an approximate monotonicity constraint and resets its context on encountering new lines, as described in~\autoref{section:entropy-context}. All models are allocated the same computational (flop) budget during training. For the {\textsc BLT}-Entropy variant, we further tune the entropy threshold at inference, lowering it from 0.6 to 0.1, which yields improved task accuracy in exchange for an increased number of inference steps.

The {\textsc BLT}-Entropy model achieves superior performance compared to the Llama 3 model across 4 of the 7 evaluated tasks, despite both models being trained with an equal amount of data in terms of bytes. This enhancement likely stems from two primary factors: (1) more efficient utilization of training compute facilitated by dynamic patching, and (2) the explicit modeling of data at the byte level rather than at the token level.

Conversely, {\textsc BLT}-Space shows inferior performance to the Llama 3 tokenizer on all tasks except one; however, its increased average patch size of 6 bytes results in a substantial decrease in inference flops. By contrast, both bpe and entropy-patching approaches yield an average patch size of around 4.5 bytes when evaluated on the same training corpus. Given an equal training computational budget, the model utilizing a larger patch size is able to process 30\% more data relative to the other two models, which may, in turn, cause {\textsc BLT} to deviate further from the compute-optimal regime.

\subsection{Patches Scale Better Than Tokens} In {\textsc BLT} models, it is possible to scale up both the patch size and the overall model size, all while preserving an unchanged compute budget for both training and inference, as well as a fixed training dataset size. This capability to increase patch size to any desired extent is a distinctive advantage of patch-based architectures, setting them apart from traditional fixed-vocabulary, token-based approaches that inherently involve efficiency compromises (see Section~\ref{section:bpe-patch}). Employing longer patch sizes reduces the computational overhead, thus freeing resources that can be redirected to expand the global latent transformer, given that it operates with reduced frequency.

We perform a fixed inference scaling experiment to evaluate whether increasing model size, while reducing the number of steps on larger patches, yields superior performance compared to using smaller models with more steps. Beginning with models containing 400 million and 3.6 billion parameters—both utilizing the Llama 2 tokenizer—we identify computationally equivalent counterparts using the Llama 3 tokenizer, as well as {\textsc BLT}-Entropy models featuring mean patch sizes of 6 and 8 bytes across the training datamix (refer to~\autoref{tab:fixed-inf-params} for specific model configurations). For those models employing patch size 8, we replace 1 encoder layer with 3. All models are trained across a range of predetermined training flop allocations.

As illustrated in \autoref{fig:fixed_inference_scaling}, {\textsc BLT} architectures demonstrate more favorable scaling behaviors compared to traditional tokenization-based models across both categories of inference computational cost. While BPE-based models initially outperform with lower training expenditures, {\textsc BLT} quickly overtakes them beyond the compute-optimal point. From a practical standpoint, allocating greater resources for a one-off pretraining phase can yield superior models under a constant inference constraint. This phenomenon is evident in the 8B parameter class, exemplified by Llama 3.1, which was exposed to two orders of magnitude more training data than is considered compute-optimal for its scale.

The threshold at which {\textsc BLT} surpasses token-based models has shifted slightly nearer to the compute-optimal configuration with the transition to higher flop class models—moving from 3x to approximately 2.5x the optimal compute budget. Additionally, the model employing patch size 8 within this larger flop class exhibits a more rapid improvement in scaling, overtaking other models earlier. As elaborated in Section~\ref{section:parameter-matched-scaling}, increasing patch size tends to align {\textsc BLT}'s performance more closely with that of BPE models as model size increases. We partly ascribe this effect to the reduced proportion of overall flops consumed by the byte-level Encoder and Decoder modules, which appear to scale more slowly compared to the Latent Transformer. Notably, as the total number of parameters increases twentyfold—expanding from 400M to 8B—the local model parameters in {\textsc BLT} are only increased by roughly a factor of two. This observation is crucial, since expanding patch size only impacts the flops of the patch Latent Transformer, leaving the byte-level modules unaffected. This dynamic also clarifies why the ratio of {\textsc BLT}-Entropy ps=8 to Llama 2 model size rose from 1.6x to 1.7x at higher model scales.

In conclusion, our investigation into patch-length scaling reveals that the {\textsc BLT} patch-based model exhibits superior scaling behavior when both the patch dimension and the model capacity are enlarged together. These advantageous scaling patterns not only remain but appear to further strengthen as the model size increases.

\section{Byte Modeling Improves Robustness} We further evaluate the robustness characteristics of {\textsc BLT} in comparison to token-based models that do not incorporate byte-level features, and introduce a technique for adapting pretrained token-based models to handle byte inputs.
\label{sec:robustness}

\subsection{Character-Level Tasks}

One of the initial incentives for developing byte-level models stemmed from their inherent resilience to noise present at the byte level in input data, along with their capacity to recognize the subcomponents of tokens—an area where conventional tokenizer-based approaches often falter. To systematically assess these capabilities, we conduct supplementary experiments using benchmark datasets designed to test robustness against input noise and sensitivity to characters across languages, such as English and other multilingual scripts, as well as including numerals and phonetic elements. The outcomes of these evaluations are summarized in Table~\ref{tab:char_tasks}.

\paragraph{Noisy Data} To evaluate robustness, we generate perturbed versions of the benchmark classification datasets outlined in Section~\ref{section:evals}, directly contrasting the performance of tokenizer-based models and {\textsc BLT}. Five character-level noise techniques are implemented to systematically alter the input text: (a) \textit{AntSpeak}: The text is rendered in uppercase and each character is separated by a space. (b) \textit{Drop}: Ten percent of the text's characters are randomly eliminated. (c) \textit{RandomCase}: Randomly assigns uppercase or lowercase to half of the characters, thereby introducing case inconsistency. (d) \textit{Repeat}: For twenty percent of the characters, each is duplicated up to a maximum of four times. (e) \textit{UpperCase}: Converts every character to uppercase. When assessing the models, each noise method is individually applied to the prompt, the completion, or both, and the results are averaged across these conditions. Results summarized in Table \ref{tab:char_tasks} for the noised HellaSwag~\citep{Eval_hellaswag} task reveal that {\textsc BLT} consistently demonstrates greater resilience than tokenizer-based counterparts, yielding an average improvement of 8 points over models trained with identical data, and even surpasses the performance of the Llama 3.1 model that leverages a considerably larger training set.

\paragraph{Phonology - Grapheme-to-Phoneme (G2P)} We evaluate {\textsc BLT}'s proficiency in converting sequences of graphemes—i.e., the individual characters of a word—into their corresponding phonemic representations. Table \ref{tab:char_tasks} displays performance on the G2P task using five-shot examples, as measured by Phonology Bench~\citep{suvarna2024phonologybench}. The results indicate that {\textsc BLT} achieves superior performance relative to the baseline Llama 3 1T tokenizer-based model.

\paragraph{CUTE}
To evaluate character-level reasoning, we test {\textsc BLT} on the CUTE benchmark~\citep{edman2024cute}, which consists of multiple tasks grouped into three main types: analyzing compositionality, recognizing orthographic similarity, and manipulating character sequences. This suite of tasks is particularly demanding for models based on tokenization, as these models tend to know the spelling of their tokens but often fail to harness this information for text manipulation purposes. As shown in Table~\ref{tab:char_tasks}, {\textsc BLT}-Entropy surpasses both BPE-based Llama 3 variants by over 25 points on this benchmark. Notably, our model excels in manipulation tasks at the character level, attaining 99.9\% accuracy on both spelling benchmarks. This considerable performance gain—despite {\textsc BLT} being trained on just one-sixteenth the data compared to Llama 3.1—suggests that learning fine-grained character details is particularly challenging for BPE-based approaches. Figure \ref{fig:cute_examples} provides examples in which the Llama 3 tokenizer model fails, yet our {\textsc BLT} model handles them successfully. Word deletion and insertion are the only categories where BPE approaches yield superior results; however, while byte-level models may find word-level manipulations less intuitive, the performance difference is moderate, implying that building compositional understanding from characters to words might be a more natural progression than starting from words down to characters. Our evaluations for each task employ the same setup and use the original Huggingface prompts. Note that BPE-based models could potentially benefit from further prompt engineering.

\paragraph{Low Resource Machine Translation} We assess the performance of {\textsc BLT} on translation tasks involving both directions for twenty-one low-resource languages across six widely studied language families, using a range of writing systems from the FLORES-101 benchmark~\citep{goyal2022flores}. SentencePiece BLEU scores are reported in Table~\ref{tab:flores}. The findings indicate that {\textsc BLT} surpasses a comparable model utilizing the Llama 3 tokenizer, yielding a 2-point improvement for translations into English and a 0.5-point gain for translations from English. While performance on high-resource language pairs is on par with or marginally exceeds that of Llama 3, {\textsc BLT} achieves notably better results for many low-resource language pairs. These outcomes highlight the benefits of byte-based modeling in capturing long-tail byte sequence generalization across underrepresented languages.

\subsection{Training {\textsc BLT} from Llama 3}
\label{sec:distillation}
We investigate a methodology in which {\textsc BLT} models benefit from the weights of existing pre-trained tokenizer-based architectures, facilitating both faster and more stable convergence during training. To implement this, we initialize the global transformer parameters in {\textsc BLT} with those from a pre-trained Llama 3.1 checkpoint. The global transformer's parameters are then fine-tuned at one-tenth the learning rate utilized for training the local encoder and decoder, using the optimal number of steps determined for Llama 3. Table~\ref{tab:distill} reports results from this approach compared against a standard {\textsc BLT} baseline. The findings indicate that {\textsc BLT} initialized from Llama 3.1 not only surpasses the performance of both Llama 3 and the baseline {\textsc BLT}, given equivalent compute budgets, but also outperforms our {\textsc BLT}-Entropy variant (see Table~\ref{tab:evals})—even though the latter was trained on a substantially larger corpus (1T tokens). This suggests that initializing {\textsc BLT} from Llama 3.1 offers a highly efficient means of reducing training compute requirements while still achieving superior results on benchmarks such as the MMLU task.

This approach can alternatively be interpreted as adapting models that rely on tokenizers into models that do not require them, thus reconfiguring a pre-trained LLaMA 3.1 model into a {\textsc BLT} model. For a thorough evaluation, we present results for the original LLaMA 3.1, trained on 15T tokens, in Table \ref{tab:distill}, and assess its performance relative to the {\textsc BLT} variant based on LLaMA 3. While our model demonstrates a modest decrease in scores on MMLU and HumanEval, it exhibits a notably greater reduction in performance on additional benchmarks. This highlights the necessity for further research focused on fully utilizing the potential of the pre-trained model, specifically through more effective tuning of data mixtures and other hyperparameters.

\section{Ablations and Discussion}
\label{section:ablations}
Here, we examine ablation studies that support the architectural decisions made for {\textsc BLT}, as well as the patching mechanism and selection of hyper-parameters employed in training the {\textsc BLT} model with 8B parameters on the {\textsc BLT}-1T dataset.

\paragraph{Entropy Model Hyper-parameters} In order to investigate how the size of the entropy model and the length of the context window influence scaling efficiency, we train a series of byte-level entropy transformer models ranging from 1 million to 100 million parameters and employ context windows between 64 and 512 tokens. We display scaling laws—curves of bpb versus training FLOPs—generated using our $400m$ and $1b$ {\textsc BLT} models trained on the Llama-2 dataset, as shown in \autoref{fig:entropy_ablation}. Our analysis shows that improvements in scaling are associated with increases in both model size and context window length, although gains tend to plateau when exceeding 50 million parameters.

\paragraph{Types of Patching}

We conduct an ablation study of the four patching methods described in Section~\ref{section:patching}, which include: 1) Strided Patching with strides of 4 and 6, 2) Whitespace-based Patching, 3) Patching via the BPE Tokenizer corresponding to the Llama 3 tokenizer, and 4) Entropy-based patching employing a compact byte-level LLM.

Although dynamic patching decreases the effective sequence length, we ensure comparable context lengths across all patching methods by controlling for sequence length. Consequently, each model is exposed to an equivalent expected number of bytes per sequence during both training and inference, thereby eliminating confounding effects related to context size. The ablation outcomes are illustrated in Figure~\ref{fig:scaling-compute-opt}. Notably, all alternative patching strategies surpass static patching, with space patching performing nearly as well as dynamic entropy-based patching.

\autoref{tab:evals_ablation} summarizes the benchmark performance of {\textsc BLT} models using three different approaches: tokenizer-based models, space patching, and entropy-based patching. All models were trained on the Llama 2 dataset for an empirically determined optimal training duration~\citep{dubey2024llama}. While space patching adopts a more straightforward methodology and does not require the additional computational overhead of real-time entropy modeling during training, our analysis demonstrates that the improvements identified for entropy-based patching in scaling experiments~(see Section~\ref{section:scaling}) are also reflected in performance on downstream tasks.\footnote{Space patching outcomes were obtained from initial experiments without cross-attention; comparable results persist when cross-attention is incorporated.}

\paragraph{Cross-Attention} 
In~\autoref{tab:crossattn_ablations_1b}, we present an ablation study examining the placement of cross-attention modules at different layers within the encoder and decoder of {\textsc BLT}. Regarding encoder-side cross-attention, we investigate three initialization strategies for the queries: (1) assigning the same trainable embedding to all global states, (2) utilizing hash embeddings based on the bytes contained within each patch, and (3) applying pooling over the encoder’s hidden states for the bytes of a patch at the particular encoder layer.

Our results indicate that incorporating cross-attention within the \textit{decoder} yields the strongest performance gains. When applied to the encoder, cross-attention shows marginal benefits, and these are observed primarily when query vectors are initialized via pooling. Furthermore, the advantages of cross-attention become especially pronounced on the Common-Crawl dataset and are further amplified when employing larger patch sizes.

\paragraph{n-gram Hash Embeddings} 

We conduct an ablation study varying the n-gram hash embedding vocabularies across 0, 100K, 200K, and 400K, with corresponding outcomes presented in Table~\ref{tab:ngram_ablations_1b}. Our results demonstrate that hash embeddings offer improvements across all evaluated domains, showing notable benefits in Wikipedia and Github, where we observe a 0.04 bpb advantage relative to only a 0.01 bpb increase after 15k training steps at the 8B scale. Furthermore, adjusting the hash size from 500K to 300K at the 8B scale leads to only a marginal change of approximately 0.001 bpb after 15k steps, suggesting that hash embeddings play a crucial role in closing the performance gap between {\textsc BLT} and tokenizer-based approaches. However, performance gains plateau once the hash count surpasses 300K. Lastly, our analysis suggests that the improvements provided by hash embeddings are largely additive to those achieved via cross-attention, as these methods contribute gains on distinct datasets.

\paragraph{Local Model Hyperparameters} Table \ref{tab:local_layers} presents an ablation study on the effects of altering the number of layers in both the local encoder and decoder components. When utilizing hash n-gram embeddings, {\textsc BLT} achieves strong performance with a minimally complex encoder—consisting of only a single layer—while benefiting from a more robust, multi-layered decoder.

\section{Related Work}
\label{section:related_work}

\textbf{Character-Level RNNs:} Character Language Modeling has long been a widely-studied task, dating back to some of the earliest neural network approaches~\citep{sutskever2011generating, mikolov2012subword, graves2013generating}, primarily because these models can inherently handle out-of-vocabulary tokens without relying on back-off strategies. In addition, \cite{kim2016character} propose a system that exclusively encodes characters on the input side, utilizing convolutional and highway layers that feed into LSTM-based RNN architectures. Their approach achieved performance comparable to the RNN-based state-of-the-art English language models of that period, and even surpassed such models for languages with rich morphological structures—an appealing property of character-level large language models. Further, \cite{kenter2018byte} demonstrate that byte-level LSTM models for machine comprehension tasks can outperform word-based counterparts on highly inflected languages such as Turkish and Russian. In related work, \cite{zhang2015character} employ character-level convolutional networks for classification problems, reporting superior results over word-level models for specific tasks. The study by \citet{chung2019hierarchical} introduces hierarchical LSTM frameworks with boundary detectors at each stratum, enabling the discovery of latent hierarchical structures in text and yielding advancements in character-level language modeling performance. Diverging from attention-based architectures, ByteNet \cite{kalchbrenner2016neural} adopts convolutional neural network layers at the character level for machine translation tasks.

\textbf{Character-Level Transformers:} The introduction of attention mechanisms in transformer architectures~\citep{vaswani2017attention}, alongside subword tokenization techniques~\citep{sennrich-etal-2016-neural}, led to considerable advancements in neural language models' capabilities on both modeling tasks and standardized evaluations. Nonetheless, the reliance on word or subword units imposes an inherent inductive bias concerning the granularity at which these models process language. Seeking to bridge transformer innovations with earlier progress in character-based modeling, \citet{al2019character} employed highly deep transformer models, leveraging auxiliary loss functions during training. Their approach enabled transformer-based character language models to achieve superior results compared to earlier LSTM-based systems. Despite these improvements, a notable performance gap persisted in relation to word-level LLMs. Additionally, GPT-2~\citep{radfordgpt2} found that, for extensive datasets such as the 1 billion word benchmark, language models operating at the byte level did not match the efficacy of their word-level counterparts.

\cite{Choe2019BridgingTG} showed that transformer-based large language models (LLMs) operating at the byte level can surpass subword-level LLMs with similar parameter counts, but these models demand significantly greater computational resources and training time. In a similar vein, \cite{el2020characterbert} trained CharFormer—a BERT-style model that constructs word-level representations by convolving over character embeddings—and achieved gains within the medical domain, albeit at a notable computational cost. CANINE, introduced by \cite{clark2022canine}, is an encoder-only model with 150 million parameters that processes raw character sequences. The architecture includes a deep transformer core—analogous to our global component—as well as a local transformer module and strided convolutions to reduce the input dimensionality. CANINE outperforms its token-level, encoder-only counterpart, mBERT, on downstream multilingual evaluation tasks. Byte-level encoder-decoder models, as investigated by ByT5~\citep{xue2022byt5}, forego patching mechanisms, resulting in improved robustness to input noise and competitive performance with tokenizer-dependent models using one-quarter of the training data. However, omitting patching requires attention computations for each byte, leading to substantial computational overhead. Modeling input at the byte granularity increases sequence lengths, which complicates efficient scaling. Recently, with the advent of the Mamba Architecture~\citep{gu2023mamba}—capable of sustaining a fixed-size memory state over extended contexts—\cite{wang2024mambabyte} trained a byte-level Mamba model without applying patching, surpassing byte-level transformer baselines (in bits-per-byte) at 350 million parameters on several benchmarks under matched compute budgets.

\textbf{Patching-based approaches:} Leveraging patching techniques can significantly reduce the high computational cost, measured in flops, associated with byte-level LLMs without necessarily sacrificing accuracy. Several studies have reported promising results at relatively modest model scales and limited training data volumes. For example, \cite{nawrot-etal-2022-hierarchical} explore the use of static patching strategies—specifically, downsampling and upsampling—in their proposed hourglass transformer architecture, which surpasses other byte-level models at the 150M parameter level. Building on this foundation, \cite{nawrot-etal-2023-efficient} introduce dynamic patching mechanisms, which include an end-to-end trainable boundary-predictor, versions supervised by predefined tokenizers, as well as an entropy-based patching model in the vein of {\textsc BLT}. Their findings suggest these dynamic approaches outperform contemporary vanilla transformer architectures in language modeling tasks at the 40M parameter scale with approximately 400M tokens of training data. Additionally, \cite{lester2024training} explore the use of arithmetic coding for sequence compression, attaining greater compression rates than possible with BPE. Through the employment of an equal-info windows method, they demonstrate better performance than byte-level baselines for language modeling, although their approach does not yet surpass results from subword-level baselines.

This research takes significant motivation from MegaByte~\citep{yu2023megabyte}, a purely decoder-based causal large language model that employs a predetermined, unchanging scheme for patching and concatenating representations to transition from bytes to patches. On the decoder side, MegaByte utilizes a localized model to revert patches to their original byte form. Their results indicate that MegaByte achieves parity with models utilizing tokenizers at the 1B parameter level on datasets encompassing approximately 400B bytes. In our analysis, we perform an ablation of MegaByte across all our experiments and observe that its use of static patching does not reach the performance of contemporary, compute-efficient models based on tokenizers under a fixed FLOPs budget. Through our approach with {\textsc BLT}, we illustrate a means to overcome this limitation. Parallel findings are reported by \cite{slagle2024spacebyte}, who also highlight the limitations of MegaByte and propose augmenting static patching by targeting whitespaces and similar space-indicative bytes, in addition to incorporating a local encoder. Their empirical results demonstrate gains over transformer models trained on tokenized inputs, particularly within domains like Github and arXiv at the 1B scale under constrained computation. In this work, we provide further experiments with this model and show that advancing the architecture remains critical in order for byte-level models to be fully competitive with the present leading token-based architectures such as Llama 3.

\section{Limitations and Future Work}

In this study, we select model training durations based on the optimal step counts established for Llama~3~\citep{dubey2024llama} to guide architectural decisions. Nevertheless, these scaling laws were originally derived for BPE-level transformer architectures, which might not yield the most effective (data, parameter size) ratios when applied to {\textsc BLT}. Determining scaling laws tailored to {\textsc BLT} remains a topic for future investigation, as such analyses may reveal even more advantageous scaling patterns for our proposed architecture. It is also important to note that the majority of our experiments utilize model sizes up to 1B parameters, leaving open the possibility that ideal architectural parameters could shift when expanding to 8B parameters or larger, potentially resulting in superior performance at greater scale.

Current transformer libraries and codebases are primarily optimized for tokenizer-based transformer models. Although our theoretical experiments maintain flop equivalence and we employ some optimized solutions (for example, FlexAttention) to accommodate components that differ from standard transformer layers, our current implementations might still lag behind tokenizer-based models concerning actual wall-clock performance and could see improvements with further optimization efforts.

Although {\textsc BLT} employs an independently trained entropy model for the patching process, integrating patching model training into a fully end-to-end pipeline presents a promising avenue for subsequent investigation. As detailed in Section \ref{sec:distillation}, we provide preliminary experimental results that demonstrate promising outcomes when applying the ``byte-ifying'' approach to tokenizer-based architectures like Llama 3, which have been pre-trained on datasets exceeding 10T tokens, by transferring and freezing the parameters of their global transformer component. Advancing this line of research may reveal techniques that not only preserve the advantages of bytefication, but also achieve higher performance relative to these tokenizer-based systems, circumventing the need for training from scratch.

\section{Conclusion}
\label{section:conclusion}

In this work, we introduce the Byte Latent Transformer (\textbf{BLT}), a novel architecture that challenges the traditional reliance on fixed-vocabulary tokenization within large language models. By employing an adaptive, learnable strategy to aggregate bytes into patches, {\textsc BLT} assigns computational attention according to the inherent complexity of input data. This mechanism yields notable gains in both computational efficiency and model robustness. Through comprehensive scaling experiments, we show that {\textsc BLT} achieves parity with established tokenization-based models such as Llama 3 at configurations up to 8B parameters and 4T bytes processed, while also being able to trade negligible declines in standard metrics for reductions of as much as 50\% in inference computational cost. Moreover, {\textsc BLT} facilitates a novel approach to model scaling—enabling the concurrent enlargement of both model and patch sizes under a fixed inference budget—offering particular benefits for compute-constrained settings prevalent in real-world deployments. By directly processing byte-level inputs, {\textsc BLT} further enhances the model's handling of infrequent or noisy data, boosting resilience to non-standard or corrupt input and supporting a more nuanced capture of sub-word information. Collectively, these findings position {\textsc BLT} as a strong candidate to replace traditional tokenization-based methods, offering an adaptable and scalable foundation for more robust and efficient language modeling.

\end{document}